\chapter{Measurement of anomalous couplings in top-quark pair decays} \label{ch::Analysis}

The previous chapter has described in detail the technicalities of the Matrix Element method and demonstrated its applicability using a simple example.
The main idea behind this method is that it combines all the individual event probabilities into an overall likelihood $\mathcal{L}_{MEM}$.
This likelihood distribution is then transformed into a negative log likelihood, which is minimized in order to extract the Matrix Element estimator $\hat{\epsilon}_{MEM}$.
In this chapter, which will focus on the measurement of the right-handed tensor coupling of the Wtb interaction vertex, the estimator of interest is denoted as $\gREst$.
\\
%***********************
% Repeat MEM likelihood equation??
%***********************

Before the actual measurement can be performed, a number of tests have to be carried out in order to ensure the Matrix Element estimator is behaving properly.
Investigating the performance of this estimator, explained meticulously in Section~\ref{sec::EstimatorProp}, is crucial since the Matrix Element method can be influenced by the introduced simplifications/assumptions.
%the sometimes simplified \textit{representations}.
Determining the correspondence between the result obtained from the Matrix Element estimator and the expected $\gR$-value will be done using simulated samples, for which the Standard Model configuration should be recovered.
%Hence, comparing the obtained result of simulated samples, where the outcome is known, with the value of the estimator allows to fix the correspondence between the estimator and the actual outcome.
%
%In the ideal case, the estimator $\hat{g_{R}}$ should be identical to the $\gR$ value extracted from the event-likelihood. 
%Nevertheless, this perfect behaviour can be influenced by ... (\textit{What can be responsible for a slope different from 1? Is it also just a bias?}) and result in a ... or a bias.
%
\\
Once all the ... the measurement itself will be discussed in Section~\ref{sec::gRMeas}.

\section{Performance of the Matrix Element estimator} \label{sec::EstimatorProp} %Is method needed here? --> Otherwise call it 'Performance of the estimator $\hat{g_{R}}$'

The applicability of the Matrix Element method can be seriously hampered by either the constructed model and the approximations made herein or otherwise by the procedure developed to extract the estimator from the event likelihood.
Here the two possible deviations from the perfect behaviour will be discussed separately, first the linearity test which focusses on the slope and afterwards the presence of an offset will be checked.

\begin{equation}
 \hat{g_R} = a \cdot \gR^{Lik} + bias
\end{equation}
The final performance check will focus on the statistical properties of the estimator.

\subsection{Linearity test}

%Why separate linearity test?
The choice to perform the linearity test separately from the offset determination is driven by fact that the linearity test requires information from several samples generated with different coupling coefficients. Since creating such samples for reconstructed events is a non-trivial task, especially compared to the ease with which generator-level samples can be created, it has been opted for in this thesis to use generator-level events for the linearity test.
The influence of the reconstructed events will be incorporated afterwards in the offset determination, which only requires the result of the Standard Model configuration.
\\

%What does the linearity test do, why a non-linear result can be obtained
The goal of the linearity test is to determine whether the outcome of the Matrix Element method depends on the coupling coefficient considered.
Such a behaviour can be expected in case the created model does not perfectly describes the physics processes taking place when moving away from the well-established Standard Model configuration.
%depicts the physics processes in a too simplified manner.
Therefore various samples will be created, each generated with a different coupling coefficient, and the obtained minimum of each sample will be compared with the value used to create it.
\\
\\
%Want to include some event selection cuts in order to take into account some discrepancies from here.
Besides a simplified model, the event-selection procedure is also a viable candidate to cause a dependency of the outcome on the coupling coefficient.
This because the selection efficiency might change significantly due to the altered kinematics of the Wtb interaction when varying the coupling coefficients.
However, since generator-level events are considered the full event-selection chain cannot be applied but, in order to mimic the conditions of the reconstructed events partially, the generated events will have to fullfill some basic event-selection criteria. These have been summarised in Table~\ref{table::GenCuts}.
The identification and additional fine-tuning criteria cannot be taken into account, but are on the other hand not expected to depend on the considered value of the coupling coefficient.
\begin{table}[h!t]
 \centering
 \caption{Basic event selection applied to the generator-level events in order to partially mimic the situation existing for the reconstructed collision events.} \label{table::GenCuts}
 \renewcommand{\arraystretch}{1.2}
 \begin{tabular}{c|c|c|c}
  Particle 	& $\pT$-cut 		& $\vert \eta \vert$-cut 	& $\Delta$R-cut 	\\
  \hline
  Jets 		& $>$ 30 $\GeV$ 	& $<$ 2.5			& $<$ 0.3		\\
  Muon		& $>$ 26 $\GeV$		& $<$ 2.1			& $<$ 0.3		\\
  Neutrino 	& $>$ 25 $\GeV$		& $<$ 2.5			& $<$ 0.3		
 \end{tabular}
\end{table}

%What is done and what is expected
The generated samples, each containing 20 000 events in order to be comparable in size to the data sample, are then analysed by the Matrix Element method and the correspondence between the obtained minimum and the coupling coefficient imposed during the generation process is determined. 
In the ideal case these two values should be identical, implying that both the model and the procedure developed to extract the estimator are behaving properly.
The outcome of the linearity test for this analysis is given in Figure~\ref{fig::CalibCurve}, which clearly indicates that the distribution of the minima can be described by a straight line with slope almost equal to $1$ ($a$ = 0.9713).
\\
%What is shown in the figure, and which range is important.
The fit on the distribution, the actual linearity test, has only been applied in the range $\left[-0.17, 0.17\right]$ since this is the relevant region to where the Standard Model configuration should be recovered. However, due to the precision of the obtained results, the region of interest can be limited further to values of $\vert \gR \vert$ smaller than 0.1, for which the linear behaviour is clearly visible.
\\
%Why deviation outside range!
The deviation from the expected distribution outside the considered range are most likely caused by the simplifications applied in the constructed model. In order to perfectly describe the physics processes over the entire $\gR$ range a more profound theoretical description of the Wtb interaction vertex is required, which lies beyond the scope of this thesis.
\begin{figure}[h!t]
 \centering
 \includegraphics[width = 0.65 \textwidth]{Chapters/Chapter6_Analysis/Figures/CalibrationCurve_SlopeComparision_DifferentCuts.pdf}
 \caption{Outcome of the linearity test based on generator-level events, for wich the obtained curve is perfectly described by a straight line. Hence both the developed model and method behave as expected and no calibration is required.} \label{fig::CalibCurve}
\end{figure}
%********************************
%Question: Fit also done by excluding the outer bins?
%********************************
%

%Some conclusions
The performed linearity test has clearly indicated that both the model and method are behaving accordingly and that the obtained results are independent of the considered coupling coefficient.
Hence, no calibration is required to restore the expected linear behaviour of the Matrix Element method.

\subsection{Offset calibration}

The aim of this second performance test is to determine whether the outcome of the Matrix Element estimator is \textit{shifted with an offset}.
Even though the linearity test has proven that an excellent agreement exists when analysing generator-level events, both for the slope and the bias, this cannot be generalised to reconstructed events.
Since this bias might be affected by the different nature of reconstructed events, its value should be determined using the simulated events fullfilling the entire event-selection chain introduced in Chapter~\ref{ch::EvtSel}.
\\
Applying the full analysis procedure on this sample of recontstructed events would ideally result in the Standard Model configuration to be retrieved.
In case a deviation from the expectation is observed, this difference will be taken into account as a bias of the developed model and the final result will be calibrated accordingly. 
\\

However, the obtained result was inconsistent with the Standard Model prediction, implying that the developed analysis procedure should be adapted in order to correctly handle reconstructed events.
Since the overall $\NegLL$ distribution corresponded to a decreasing line with minimum located at the edge of the considered range, the observed discrepancy is more than merely a bias introduced by the \textit{procedure}.
\\
A closer look at the issue indicated that the discrepancy is most likely caused by a certain type of events the Matrix Element method is unable to analyse properly.
Hence before the actual measurement can proceed/be performed, the reason behind this deviating/specific behaviour should be understood thoroughly such that a procedure can be developed to limit their influence.

\paragraph{(Specific) characteristics of reconstructed events} \hfill \\ %Understanding the nature of the \textit{bad} events

The inconsistent outcome obtained for generator-level and reconstructed events suggest that the Matrix Element method behaves differently for the two types of events.
This is not completely unexpected since reconstructed events tend to be influenced by detector inefficiencies or ill-determined event kinematics.
However, the Matrix Element will treat all events as actual semi-leptonic top-quark pair decays, and any deviation from the expected topology is thus likely to result in an incorrect event probability.
\\
The performed study will first focus on a sample of selected $\ttbar$ events for which the four jets have been correctly matched with the generator-level parton to exclude the contribution from badly reconstructed event topologies.
In case the bias is also observed for this sample, (it) can be concluded that the deviation is definitely inherent to the nature of reconstructed events and not caused by a lower reconstruction efficiency.
\\
\\
In order to determine whether the observed discrepancy is caused by a specific type of events with a distinct signature, an exhaustive study has been performed with the aim of finding an explanation for the significant disagreement between the generator- and reconstructed-level measurement.
The most promising variable encountered to differentiate between well and less behaving events is the value of the $\NegLLEvt$ distribution evaluated for each event at the Standard Model configuration.
\\
This distribution has a significantly different shape for generator-level and reconstructed events, as can be seen in Figure~\ref{fig::SMLik}.
For the correctly reconstructed $\ttbar$ events, shown on the left, a tail is clearly visible, which is on the other hand completely absent for the generator-level distribution on the right.
The reconstructed-level distribution seems to suggest that a fraction of events have a lower event probability, and thus higher $\NegLLEvt$ value, assigned in case the Matrix Element is unable to correctly analyse the kinematics of the considered event.
%***************************
% Necessary to mention something about the different starting value??
%***************************
\\
\begin{figure}[h!t]
 \centering
 \includegraphics[width = 0.45 \textwidth]{Chapters/Chapter6_Analysis/Figures/SMLikelihoodValue_GenEventsSM.pdf} \hspace{0.3cm}
 %Taken from directory: Events_CalibCurve/CalibCurve_SemiMu_RgR_AllDeltaTF_MGSampleSM_20000Evts_CutsAlsoOnMET/SMLikelihoodValue_GenEventsSM.pdf
 \includegraphics[width = 0.45 \textwidth]{Chapters/Chapter6_Analysis/Figures/SMLikelihoodValue_RecoEvents_MyTF.pdf}
 %--> Update this!!
 \caption{Distribution of the $\NegLLEvt$ value for the $\gR$ = $0.0$ configuration for both the generator- (left) and reconstructed-level (right) events.} \label{fig::SMLik}
\end{figure}

In order to ensure this behaviour is not influenced by wrongly reconstructed event topologies, the same procedure has also been applied onto a sample containing simulated $\ttbar$ events for which at least one jet has not been matched with the correct generator-level parton.
This allowed to compare both the outcome of the Matrix Element estimator and the distribution of this $\NegLLEvt$ variable with the previously considered sample containing only correctly reconstructed $\ttbar$ events, which were both identical.
Hence, it can be concluded that the incorrect determination of the event topology is not causing the peculiar behaviour of the Matrix Element method, but that it is an undesirable drawback of applying this complex technique on reconstructed collision events.
\\
Figure~\ref{fig::SMLikCorrVSWr} contains the normalised distributions of the $\NegLLEvt$ variable evaluated at the SM configuration for the two considered $\ttbar$ samples.
Even though some differences clearly exist between the two samples, the position of the peak and the \textit{importance/relevance/size} of the intermediate region, the tail is almost identical.
So this comparison demonstrates that the Matrix Element method definitely treats events with a wrong event topology differently, but also that these type of events are not the responsible for the discrepancy observed for reconstructed events.
\\
\begin{figure}[h!t]
 \centering
 \includegraphics[width = 0.7 \textwidth]{Chapters/Chapter6_Analysis/Figures/ScaledContribution_CWU_LumiNorm.pdf}
 \caption{Normalised distribution of the $\NegLL$-variable evaluated at the Standard Model configuration for correctly reconstructed (green) and wronly reconstructed $\ttbar$ events. \textit{Should still remove data and unmatched! + SCALE TO 1!}} \label{fig::SMLikCorrVSWr}
\end{figure}

A final study which has been performed attempts to find an explanation for the peculiar behaviour observed for the reconstructed events by investigating the dependence of other event variables on this $\NegLLEvt$.
The goal of this study is to search for a distinguishing feature that could be used to reject the events negatively affecting the output of the Matrix Element method.
In order to be unaffected by other influences, only the $\ttbar$ sample containing events for which the topology is correctly reconstructed has been considered.
\\
It has been decided to focus on the distribution of the $\NegLLEvt$ variable, which sould correspond in the ideal case to a parabola with minimum value located at $\gR$ = 0.
Hence the different variables which have been considered are related to the steepness of this distribution such that can be visualised for how many events the expected behaviour is recovered.
The dependence of these variables on the $\NegLLEvt$ variable are shown in Figure~\ref{fig::SMLik2D}, where this variable is placed each time on the $x$-axis.
However, it should be noted that studying event-based variables for the Matrix Element method is rather challenging since this technique requires enough statistics in order to be sufficiently accurate/precise.
\\
\\
The upper left histogram contains the maximum difference observed for the $\NegLLEvt$ distribution, thus indicating the significance of this event in the overall Matrix Element method likelihood since events with a large \textit{amplitude} influence the overall shape the most. 
This clearly indicates that most events have a rather flat shape, with the exception of some outliers mainly located at low values of this $\NegLLEvt$ variable. However, a significantly deviation behaviour is visible once the $\NegLLEvt$ variable becomes larger than approximately $65$, which is also the position of the tail in Figure~\ref{fig::SMLikCorrVSWr}.
\\
The upper right histogram shows the value of the second derivative when a parabola is drawn through the point $\gR$ = $0$ and the two points surrounding it ($\gR$ = $\pm$ $0.05$).
For the main bulk of events this value is close to $0$, but still slightly positive, indicating that the corresponding parabola is rather flat. However, for higher values of this $\NegLLEvt$ variable, the value of the second derivative starts to behave unexpectedly.
\\
In the lower two histograms the difference of the $\NegLLEvt$ variable between the Standard Model configuration and the outer left or right $\gR$ value is plotted.
This value is slightly related with the maximum observed difference, with the exception that the former one depends more on the expected shape of the $\NegLLEvt$ distribution.
In order to have a nice parabola with minimum located at $\gR$ = 0, implying that this difference should be negative for the outer left $\gR$ variable and positive for the outer right one.
Also here the desired behaviour is recovered partially for events with a low $\NegLLEvt$ value, but starts to act rather abnormal once this variable becomes larger than approximately $65$.
\\
\begin{figure}[h!t]
 \centering
 \includegraphics[width = 0.45 \textwidth]{Chapters/Chapter6_Analysis/Figures/LnLikCut_ScatterPlot_SMLikvsMaxDelta_AllTT.pdf}
 \includegraphics[width = 0.45 \textwidth]{Chapters/Chapter6_Analysis/Figures/SMLik_vs_ScdDerFine_CorrectTTbar.pdf}
 \includegraphics[width = 0.45 \textwidth]{Chapters/Chapter6_Analysis/Figures/SMLik_vs_LeftDeltaLnLik_CorrectTT.pdf}
 \includegraphics[width = 0.45 \textwidth]{Chapters/Chapter6_Analysis/Figures/SMLik_vs_RightDeltaLnLik_CorrectTT.pdf}
 \caption{2D plot for all $\ttbar$ events (so CWU combined -- scdDer only for correct ...). The two lower plots contain on the y-axis the $\chiSqMEM$ difference between $\gR$ = 0.0 and $\gR$ = $\pm$ 0.2. Hence in order to have minimum at zero, both should actually be negative!} \label{fig::SMLik2D}
\end{figure}

Unfortunately, the different variables considered did not result in conclusive/sufficient(?) proof in order to decide which effect is causing the Matrix Element method to misdetermine the event probabilities of some specific events.
However, the presented distributions clearly indicated that for each of these variables some unexpected effects takes place once the $\NegLLEvt$ becomes larger than a specific value.
Hence, the only option to reduce the influence of the events the Matrix Element is incapable of handling, is by limiting the value of this $\NegLLEvt$ variable for each event.
This is necessary since the events located in the tail of the distribution become rather important because they contribute on average the most to the overall $\NegLL$ distribution from which the value of the Matrix Element estimator is \textit{extracted/determined}.


\paragraph{Matrix Element event-cleaning procedure} \hfill \\

The required event-cleaning procedure should be developed carefully since it should only exclude the events residing in the actual tail of the $\NegLLEvt$ distribution.
No attempt should be made to reduce the contribution of badly reconstructed event topologies since this is already taken care of by the event selection criteria formulated in Chapter~\ref{ch::EvtSel}.
Especially since the fraction of these type of events is almost negligible compared to the correctly reconstructed events, as can be seen in Figure~\ref{fig::SMLikCorrVSWrUnSc}. 
This shows the distribution of the $\NegLLEvt$ variable for both type of events, as was the case in Figure~\ref{fig::SMLikCorrVSWr}, without normalising the number of entries.
Hence the relative importance of the two considered samples is clearly visible, which is dominated by the sample containing the correctly reconstructed event topologies
\\
\begin{figure}[h!t]
 \centering
 \includegraphics[width = 0.65 \textwidth]{Chapters/Chapter6_Analysis/Figures/RelativeContribution_CWU_LumiNorm.pdf}
 \caption{Distribution of the $\NegLLEvt$ variable evaluated at the Standard Model configuration for correctly reconstructed (green) and wronly reconstructed $\ttbar$ events. \textit{Should still remove data and unmatched!}} \label{fig::SMLikCorrVSWrUnSc}
\end{figure}

From the performed studies could be concluded that the output of the Matrix Element method starts to behave incorrect as soon as the value of the $\NegLLEvt$ variable becomes larger than approximately $65$.
However, in order to choose the optimal cut-value to apply, it has been opted for in this thesis to determine the influence of several cut-values and select the one for which no offset exists.
For this the entire collection of selected events passing the full event-selection chain will be used.
\\
This is an acceptable approach since in case the Matrix Element method would have capable of dealing with reconstructed events from the start, any observed bias would have been corrected for.
The procedure followed now will adapt the output of the Matrix Element method in order to ensure these type of events are treated correctly and thus unaffected by an offset.
%In contrast, the procedure followed now requires to adapt the output of the Matrix Element method in order to be capable of analysing these type of events.
%Hence it is correct to choose the necessary improvements in such a way that the outcome obtained from the Matrix Element estimator $\gR$ correspond with the Standard Model prediction.
\\

The actual determination of this cut-value starts by scanning over the relevant region of the $\NegLLEvt$ distribution.
Since the previous studies have indicated that this is most likely located around $65$, the considered range has been restricted to $\left[60, 70\right]$.
\\
CONTINUE FROM HERE!\\
Hence, several cut-values have been applied onto this sample for which the corresponding minimum has been determined, as can be seen in Figure~\ref{fig::CutValueFit}.
The considered cut-values range between 60 and 70, and can thus all be found in the outer range of the $\chiSqMEM$ distribution.
\\
Figure~\ref{fig::CutValueFit} also indicates the strong dependence of the obtained minimum on the applied cut-value, implying that its optimal value should be determined with great care in order to avoid introducing a significant bias. Hence, the events located more in the intermediate region of the $\chiSqMEM$ distribution, between $60$ and $65$ approximately, still contain relevant information and should be preserved for further analysis.
\\
\begin{figure}[h!t]
 \centering
 \includegraphics[width = 0.65 \textwidth]{Chapters/Chapter6_Analysis/Figures/MinComp_MCOnly_StraightLineAtZero.pdf}  %Data lumi used (be consistent)!
 \caption{Minima obtained by applying different cut-values. This distribution has been fitted with a polynomial in order to determine the optimal cut-value, which correspond to the value for which no bias is recovered.} \label{fig::CutValueFit}
\end{figure}

The overall distribution containing the minima corresponding to each cut-value is then fitted with a polynomial function, in Figure~\ref{fig::CutValueFit}.
From this the optimal cut-value can easily be determined as being $63.869$ for which a minimum equal to $\gR$ = -0.001 $\pm$ 0.009 is recovered, which is clearly statistically compatible with 0.

%The effect of this event-cleaning procedure on the different samples has been summarised in Table~\ref{table::CutInfl}, which varies significantly for the considered samples.
%\begin{table}[h!t]
% \centering
% \caption{Influence of the event-cleaning procedure on the different samples. RELEVANT??} \label{table::CutInfl}
% \renewcommand{\arraystretch}{1.2}
% \begin{tabular}{c|c}
%  Sample 		& Event reduction 	\\
%  \hline
%  $\ttbar$ (good)	& $\%$ 			\\
%  $\ttbar$ (wrong) 	& $\%$ 			\\
% \end{tabular}
%\end{table}

%*************************
% Decide: Interesting to keep this TF part???
%*************************
%Relevant to keep?? --> Maybe move this to the conclusion part of the section??
%Next it has been investigated how such a significant tail can arise and why the Matrix Element seems to be incapable of handling a specific type of reconstructed events.
%One of the more obvious differences between generator-level and reconstructed events is the tendancy to be influenced by detector inefficiencies and ill-determined kinematic variables. For the reconstructed events this is a significant point of concern, amplified by the fact that the resolution functions of these events allow the kinematics to \textit{vary} in a much wider range.
%
%The importance of the applied resolution function can be visualised in Figure~\ref{fig::SMLikTF} which shows this $\chiSqMEM$ variable in case the resolution function developed in Section~\ref{sec::TF} is considered and otherwise in case the basic Gaussian function is used to describe the smearing of the \textbf{what exactly?}.
%Comparing these two shows a clear dependence on this resolution function, again indicating that the events in the end of the tail cannot converge since the reconstructed kinematic information does not agree with the expectation within the range allowed by the resolution function.
%\\
%\begin{figure}[h!t]
% \centering
% \includegraphics[width = 0.35 \textwidth]{image.png} %Maybe show the two on top of each other? This way repeating the same figure can be avoided!!
% \caption{Distribution of the $\chi^{2}_{MEM}$-value for the $\gR$ = $0.0$ configuration for the resolution functions created specifically for this analysis (green) and for the basic Gaussian resolution %function of the Matrix Element method (blue).} \label{Fig::SMLikTF}
%\end{figure}

%Possible conclusion:
Even though the process responsible for the deviating behaviour is not completely mastered, it seems more than plausible that the reconstructed events are heavily influenced by inefficiencies non-existing for generator-level events. Moreover, the Matrix Element method is developed to treat every event as if it is a perfectly described semileptonic top-quark pair decay. Any inconsistency from the expected topology is likely to result in the event probability being misdetermined since the mathematical framework of the method is unable to converge.
Nonetheless, it has been established that the origin of this abnormal behaviour is not caused by the applied event-selection procedure but is a \textbf{true} feature of the Matrix Element method such that a strategy should be developed to reduce the contribution of these type of events.

\paragraph{Data-MC discrepancy}
Nice agreement between data and MC (within 1 sigma interval if unweighted MC uncertainty is used). 
Shape is also almost identical for the $\chi^{2}$ distribution evaluated at the Standard Model configuration.\\
Optimal cut-value differs slightly: 63.396 for data and 63.8764 for all simulated samples.

\begin{figure}[h!t]
 \centering
 \includegraphics[width = 0.45 \textwidth]{Chapters/Chapter6_Analysis/Figures/StackedPlot_SMLikValue_AllMCAndData.pdf}
 \includegraphics[width = 0.45 \textwidth]{Chapters/Chapter6_Analysis/Figures/MinComp_MSPlot_StraightLineAtZero.pdf}            %Same lumi used (data)!
 \caption{Distribution of the $\NegLLEvt$ variable evaluated at the SM configruation point and the obtained minimum-values for the different cut-values considered both for data (red) and all simulated (black) samples.}
\end{figure}

\subsection{Statistical properties}
The developed technique to extract the relevant information from the Matrix Element method has proven to be uninfluenced by any bias, as has been shown in Section~\ref{sec::CalibCurve}.
Furthermore, the optimal cut-value for the event-cleaning procedure is chosen as such that the expected Standard Model value is recovered for the simulated signal events, top-quark pairs decay(ing semi-leptonical). Since the remaining simulation samples are all almost negligible compared to this one, the leading background contribution in this analysis is single-top production in the tW-channel which is about 50 times smaller, no significant influence on the overall outcome is expected from these types of events.
\\

However, additional research is still required in order to make sure the statistical properties of this analysis are well described.
This can be performed using a resampling technique which generates a set of samples by randomly selecting events from the considered simulated samples in accordance with a previously specified luminosity, in general the luminosity of the data sample.
This allows to obtain a representation of actual data and to verify whether the average result can be compared with expectation.
%Each of these generated samples is then treated as an actual data sample and the full analysis is applied in order to determine the distribution of both the 
\\
\\
Width of pull = 1.00395 $\pm$ 0.0240723 (using only ttbar semileptonic events)\\
Width of pull = 0.957432 $\pm$ 0.0238748 (using all ttbar events)\\
Width of pull = 0.987126 $\pm$ 0.0221517 (all MC)

\begin{figure}[h!t]
 \centering
 \includegraphics[width = 0.45 \textwidth]{Chapters/Chapter6_Analysis/Figures/PseudoExp_MinimumDistr_AllTTbarSemiLept_UpdatedXS.pdf}
 \includegraphics[width = 0.45 \textwidth]{Chapters/Chapter6_Analysis/Figures/PseudoExp_UncDistr_AllTTbarSemiLept_UpdatedXS.pdf}
 \caption{Minimum and uncertainty distribution obtained for the 1000 considered pseudo experiments. (all semiLept TT)}  \label{fig::MinAndUnc}
\end{figure}

\begin{figure}[h!t]
 \centering
 \includegraphics[width = 0.7 \textwidth]{Chapters/Chapter6_Analysis/Figures/PullDistr_AllTTbarSemiLept_UpdatedXS.pdf}
 \caption{Pull distribution} \label{fig::PullDistr}
\end{figure}

\section{Measurement of $\gR$ using the Matrix Element method} \label{sec::gRMeas}

\subsection{Data result}
\paragraph{MC results } \hfill \\

Nominal result is (Figure~\ref{fig::MinNominal}):
\begin{equation}
 \gR = -0.00139091 \pm 0.00912391
\end{equation}

\begin{figure}[h!t]
 \centering
 \includegraphics[width = 0.7 \textwidth]{Chapters/Chapter6_Analysis/Figures/FittedGraph_NominalResult_TTbarJets.pdf}
 \caption{Obtained minimum distribution for the nominal case.} \label{fig::MinNominal}
\end{figure}

Separate events on the different background samples are not always perfectly described:
\begin{table}[h!t]
 \centering
 \caption{} \label{table::BckInfl}
 \renewcommand{\arraystretch}{1.2}
 \begin{tabular}{c|c}
  Sample 						& Minimum 			\\
  \hline
  semi-leptonic $\ttbar$ 				& -0.00142159 $\pm$ 0.00912416 	\\
  semi-leptonic $\ttbar$ + full-leptonic $\ttbar$ 	&  	\\
  semi-leptonic $\ttbar$ + W-jets (4j) 			&  	\\
  semi-leptonic $\ttbar$ + Z-jets (4j) 			&  		\\
  semi-leptonic $\ttbar$ + ST tChannel 			&  	\\
  semi-leptonic $\ttbar$ + ST tW-Channel  		&  	\\
  \hline
  Total MC 			& -0.000199381 $\pm$ 0.00888694 
 \end{tabular}
\end{table}

The curve obtained when considering all MC samples is given in Figure~\ref{fig::MinNominalAll}
\begin{figure}[h!t]
 \centering
 \includegraphics[width = 0.7 \textwidth]{Chapters/Chapter6_Analysis/Figures/FittedGraph_NominalResult_AllMC.pdf}
 \caption{Obtained minimum distribution for the nominal case.} \label{fig::MinNominal}
\end{figure}

\paragraph{Actual Data result} \hfill \\

The obtained data-result is (shown in Figure~\ref{fig::MinData}):
\begin{equation}
 \gR = 0.00915763 \pm 0.00829786
\end{equation}

\begin{figure}[h!t]
 \centering
 \includegraphics[width = 0.7 \textwidth]{Chapters/Chapter6_Analysis/Figures/FittedGraph_DataResult.pdf}
 \caption{Obtained minimum distribution for the data events.} \label{fig::MinData}
\end{figure}


\subsection{Systematic uncertainties}

Obtained distribution is not perfectly Gaussian, so herefore should be corrected for.
Hence an additional systematic will be introduced which takes into account the deviation from a Gaussian distribution.
\\

Result obtained for JESMinus is (Figure~\ref{fig::MinJESMinus}):
\begin{equation}
 \gR = 0.00376009 (\pm 0.00898 \qquad \textrm{using data-lumi})
\end{equation}

\begin{figure}[h!t]
 \centering
 \includegraphics[width = 0.7 \textwidth]{Chapters/Chapter6_Analysis/Figures/FittedGraph_JESMinusResult_TTbarJets.pdf}
 \caption{Obtained minimum distribution for the JES minus case.} \label{fig::MinJESMinus}
\end{figure}

